% Einheit 5 – Kurvenuntersuchung im Sachzusammenhang
\setcounter{aufgabe}{38}
\hypertarget{e5}{}\einheitenkopf{Einheit 5 von 5 · Kurvenuntersuchung im Sachzusammenhang}
\zweigzeile{Hier lernst du, Sachfragen in Hoch-, Tief- und Wendepunkte zu übersetzen und die Ergebnisse im Sachzusammenhang zu deuten · kennst du seit Q1 (Klasse 11) · Abitur GK · baut auf: Extrempunkte (Einheit 2), Krümmung und Wendepunkte (Einheit 3), Graph und Ableitungsgraph (Einheit 4), Ableitung als Änderungsrate deuten, Gleichungen lösen, Funktionswerte berechnen}

\begin{aufgabe}{Ich erkenne, was eine Sachfrage mathematisch bedeutet. Schreibe A (Hoch- oder Tiefpunkt von $f$, Ränder beachten), B (Wendepunkt von $f$, also Extremstelle von $f'$), C ($f'(t) = 0$) oder D (Wert von $f'(t)$ berechnen). Rechne nichts.}
\begin{teile}
\teil „Wann ist die Temperatur am höchsten?“ \leerfeld
\teil „Wann nimmt die Zahl der Erkrankten am stärksten zu?“ \leerfeld
\teil „Wann ändert sich der Wasserstand gerade nicht?“ \leerfeld
\teil „Wie schnell wächst der Bestand nach 3 Tagen?“ \leerfeld
\teil „Wann nimmt die Geschwindigkeit am stärksten ab?“ \leerfeld
\teil „Welches ist der größte Wert im Zeitraum von 0 bis 10 Stunden?“ \leerfeld
\end{teile}
\end{aufgabe}

\begin{aufgabe}{Ich kann den Verlauf eines Graphen im Sachzusammenhang beschreiben. Die Abbildung zeigt die Wassermenge $W(t)$ in einem Hochbehälter (in 100 m³), $t$ Stunden nach 6 Uhr.}
Beispiel: Ein Graph steigt und wird dabei immer steiler. → Die Größe nimmt zu, und zwar immer schneller.

\begin{ksys}[xmin=0,xmax=8,ymin=0,ymax=11,ablesen,xlabel=$t$ in h,ylabel=$W(t)$ in 100 m³]
\funktionab{0.25*\x^3-3*\x^2+9*\x+2}{W}{0}{8}
\end{ksys}

Beschreibe den Verlauf im Zeitabschnitt mit „nimmt zu“ oder „nimmt ab“ und „immer schneller“ oder „immer langsamer“.
\begin{teile}
\teil $0 \le t \le 2$: \leerfeld \leerfeld
\teil $2 \le t \le 4$: \leerfeld \leerfeld
\teil $4 \le t \le 6$: \leerfeld \leerfeld
\teil $6 \le t \le 8$: \leerfeld \leerfeld
\teil Gib an, zu welcher Uhrzeit die Wassermenge am stärksten abnimmt. \leerfeld
\teil Lies ab, wie viel Wasser zu diesem Zeitpunkt im Behälter ist. \leerfeld
\end{teile}
\end{aufgabe}

\begin{aufgabe}{Ich kann größte Werte und stärkste Änderungen im Sachzusammenhang berechnen. Antworte mit Zahl, Einheit und Bedeutung.}
\beispiel{h(t) &= -5t^2 + 20t && \text{Höhe eines Balls in m, } t \text{ in s} \\ h'(t) &= -10t + 20 = 0 && \Leftrightarrow t = 2 \\ h(2) &= 20 && \text{„Nach 2 s erreicht der Ball mit 20 m seine größte Höhe.“}}
\begin{teile}
\teil Die Flughöhe einer Drohne wird für $0 \le t \le 6$ durch $h(t) = -0{,}5t^3 + 3t^2$ beschrieben ($h$ in m, $t$ in s). Berechne die größte Flughöhe.
\teil Die Zahl der Nutzer einer App wird für $0 \le t \le 14$ durch $N(t) = -0{,}1t^3 + 1{,}5t^2 + 20$ beschrieben ($N$ in Tausend, $t$ in Monaten). Weise nach, dass die Nutzerzahl nach 10 Monaten am größten ist, und gib die größte Nutzerzahl an.
\teil Die Konzentration eines Wirkstoffs im Blut wird durch $K(t) = 50t\cdot\mathrm{e}^{-0{,}5t}$ beschrieben ($K$ in mg/l, $t$ in Stunden). Berechne, wann die Konzentration am stärksten abnimmt, und gib die Abnahmerate zu diesem Zeitpunkt an.
\teil Der Wasserstand eines Flusses an einem Pegel wird für $0 \le t \le 5$ durch $f(t) = 0{,}5t^3 - 6t^2 + 18t + 20$ beschrieben ($f$ in cm, $t$ in Stunden). Bestimme den höchsten und den niedrigsten Wasserstand in diesem Zeitraum.
\end{teile}
\end{aufgabe}

\begin{aufgabe}{Ich kann größte Werte und stärkste Änderungen im Sachzusammenhang berechnen – weiter.}
\begin{teile}
\teil Das Profil eines Futtertrogs wird durch $s(x) = 0{,}1x^3 - 0{,}6x^2$ für $0 \le x \le 6$ beschrieben; der obere Rand liegt auf der $x$-Achse (1 LE = 5 cm). Prüfe, ob der Trog in einen Karton mit den Innenmaßen 32 cm Breite und 15 cm Höhe passt.
\teil Die Zuflussrate in ein Becken wird für $0 \le t \le 15$ durch $r(t) = -0{,}2t^3 + 3t^2$ beschrieben ($r$ in Litern pro Minute, $t$ in Minuten). Die Wendestelle von $r$ liegt bei $t = 5$. Berechne $r'(5)$ und deute den Wert im Sachzusammenhang. (Abitur 2025)
\teil Die Höhe zweier Pflanzen A und B wird für $0 \le t \le 12$ ($t$ in Wochen) durch $f(t) = -0{,}05t^3 + 1{,}3t^2 + t + 5$ und $g(t) = 0{,}4t^2 + 4t + 5$ beschrieben (in cm). Tom rechnet:
\rechnung{d(t) &= f(t) - g(t) = -0{,}05t^3 + 0{,}9t^2 - 3t && (1) \\ d'(t) &= 0 \;\Leftrightarrow\; t = 2 \text{ oder } t = 10 && (2) \\ d''(10) &= -1{,}2 < 0 && (3) \\ d(10) &= 10; \quad d(0) = 0; \quad d(12) = 7{,}2 && (4)}
Beschreibe die Bedeutung der Schritte im Sachzusammenhang und formuliere eine Aufgabenstellung, die Tom damit beantwortet. (Abitur 2024)
\end{teile}
\end{aufgabe}

\begin{aufgabe}{Ich finde den Fehler bei der stärksten Zunahme. Für die Temperatur $T(t) = -0{,}1t^3 + 0{,}9t^2 + 15$ (in °C, $t$ in Stunden nach 8 Uhr, $0 \le t \le 8$) sucht Paul den Zeitpunkt der stärksten Zunahme:}
\rechnung{T'(t) &= -0{,}3t^2 + 1{,}8t = 0 \\ t &= 0 \text{ oder } t = 6 \\ &\Rightarrow \text{„Um 14 Uhr nimmt die Temperatur am stärksten zu.“}}
\begin{teile}
\teil Finde den Fehler, benenne ihn und korrigiere die Rechnung.
\teil Die Zahl der Besucher eines Stadtfests wird für $0 \le t \le 9$ durch $Z(t) = -0{,}2t^3 + 1{,}8t^2 + 50$ beschrieben ($Z$ in Hundert, $t$ in Stunden nach 14 Uhr). Bestimme, wann die Besucherzahl am stärksten zunimmt, und gib die Zunahme pro Stunde an.
\end{teile}
\end{aufgabe}

\begin{aufgabe}{Ich kann Ergebnisse im Sachzusammenhang begründen.}
\begin{teile}
\teil $V(t)$ gibt das Volumen in einem Tank in Litern an, $t$ in Minuten. Begründe, warum $V'(t)$ in Litern pro Minute angegeben wird.
\teil Die Temperatur wird von 6 bis 18 Uhr gemessen. Die höchste Temperatur liegt um 18 Uhr, obwohl die Ableitung dort nicht null ist. Erkläre, wie das sein kann.
\end{teile}
\end{aufgabe}

\begin{aufgabe}{Ich kann die Maße eines umschließenden Quaders aus einem Profil bestimmen.}
\begin{teile}
\teil Der Querschnitt eines Zelts wird durch $z(x) = -0{,}5x^2 + 2$ beschrieben (1 LE = 1 m, Boden auf der $x$-Achse); das Zelt ist 3 m lang. Berechne das Volumen des kleinsten Quaders, der das Zelt umschließt.
\teil Eine Vase entsteht, wenn der Graph von $r(x) = 0{,}1x^3 - 0{,}9x^2 + 2{,}4x + 1$ für $0 \le x \le 6$ um die $x$-Achse rotiert (1 LE = 1 cm). Sie wird aus einem quaderförmigen Holzblock gedreht, der die Vase genau umschließt. Berechne die Maße des Blocks und seine Masse bei einer Dichte von 0,7 g/cm³.
\end{teile}
\end{aufgabe}

\begin{aufgabe}{Ich kann Steigung und Neigungswinkel im Sachzusammenhang berechnen. Eine Straße steigt auf 100 m Länge nach dem Profil $h(x) = -0{,}00002x^3 + 0{,}003x^2$ an ($x$ und $h$ in m, $0 \le x \le 100$).}
\begin{teile}
\teil Bestimme die Stelle, an der die Straße am steilsten ist, und gib die Steigung dort in Prozent an.
\teil Berechne den größten Neigungswinkel der Straße und prüfe, ob sie die Vorgabe „höchstens 10°“ einhält.
\end{teile}
\end{aufgabe}

\begin{aufgabe}{Ich kann Werte am Graphen im Sachzusammenhang ablesen und deuten. Links: Kosten $K(x)$ und Erlös $E(x)$ eines Betriebs ($x$ in 100 Stück, Beträge in 1000 €). Rechts: Wasserstand $p(t)$ in einem Hafen ($p$ in m, $t$ in Stunden).}

\begin{ksys}[xmin=0,xmax=9,ymin=0,ymax=60,ystep=10,ablesen,xlabel=$x$,ylabel=in 1000 €]
\funktionab{0.1*\x^3-0.5*\x^2+1.6*\x+8}{K}{0}{9}
\gerade[7]{5}{0}{E}
\end{ksys}\ksysabstand\begin{ksys}[xmin=0,xmax=24,xstep=2,ymin=0,ymax=6,ablesen,xlabel=$t$ in h,ylabel=$p(t)$ in m]
\funktionab{3+2*sin(30*\x)}{p}{0}{24}
\end{ksys}

\begin{teile}
\teil Lies ab, für welche Stückzahlen der Erlös größer ist als die Kosten. \leerfeld
\teil Deute $K(0) = 8$ im Sachzusammenhang.
\teil Lies ab, um welchen mittleren Wasserstand der Wasserstand im Hafen schwankt. \leerfeld
\end{teile}
\end{aufgabe}
